% Cote 119: offprints, reports, courses, CV and letters (typed and printed part): archive cote 119, site page 188 (PDF page 189).
% typescript. Draft compared: none.
% AI-transcribed by Claude Opus 5.5 (claude-opus-5-5, adaptive thinking, effort high), Message Batches API, batch r119_full (msgbatch_01DwZNm5Y7YsKBWyAN4sFs8w), 2026-11-13; not yet checked by eye.
\scanpage{119/188}
\mspage{-7-}

S'il existe un développement en série de Laurent celui-ci est donc unique.

Prenons donc $r'_1$ et $r'_2$ tels que $\rho_2 < r'_2 < r_2 < r_1 < r'_1 < \rho_1$ .

Soient $\gamma_1$ et $\gamma_2$ les cercles de centre 0 de rayon $r'_1$ et $r'_2$ orientés
dans le sens direct. Pour $z$ tel que $r_1 < |z| < r_2$ considérons le segment
$[r'_2\ e^{i\theta}, r'_1\ e^{i\theta}]$ pour une valeur de $\theta$ qui ne soit pas un argument de $z$.
Soit $\gamma$ le chemin fermé qui consiste à parcourir $\gamma_1$ dans le sens direct en
partant de $r'_1\ e^{i\theta}$ puis le segment $[r'_1\ e^{i\theta}, r'_2\ e^{i\theta}]$ puis $\gamma_2$ dans le sens
indirect puis le segment $[r'_2\ e^{i\theta}, r'_1\ e^{i\theta}]$. Ce chemin est homotope à un point
dans la couronne ouverte définie par $]\rho_1 , \rho_2[$. En effet, choisissons un para-%
métrage de $\gamma$ sur $[0,1]$ $\gamma(t) = \rho(t)\ e^{i\varphi(t)}$ tel que $\rho$et $\varphi$ soient continues
et $\theta \leq \varphi(t) \leq \theta+2\pi$ la fonction définie par
\[ F(t,u) = [(1-u)\ \rho(t) + ur'_1]\ e^{i(1-u)\ \varphi(t) + iu\theta} \]
est évidemment continue et $F(t,0) = \gamma(t)$ \quad $F(t,1) = r'_1\ e^{i\theta}$ .

L'indice de $\gamma$ par rapport à $z$ est visuellement 1.

On peut s'en assurer en remarquant que par définition de l'indice
on a : \qquad $I(\gamma,z) = I(\gamma_1,z) + I(\gamma_2,z) = 1+0 = 1$

les intégrales le long des segments $[r'_1\ e^{i\theta}, r'_2\ e^{i\theta}]$ et
$[r'_2\ e^{i\theta}, r'_1\ e^{i\theta}]$ s'annulant.

La proposition 6.1.1. nous donne donc :
\[ (8.2.4) \qquad f(z) = \frac{1}{2\pi i}\int_{\gamma_1} \frac{f(t)\ dt}{t-z} - \frac{1}{2\pi i}\int_{\gamma_2} \frac{f(t)\ dt}{t-z} \]
On procède comme en 5.5.

- Dans la première intégrale on a $|t| = r'_1$ et $|z| \leq r_1 < r'_1$

On a donc : $\dfrac{1}{t-z} = \sum\limits_{n \geq 0} \dfrac{z^n}{t^{n+1}}$

La convergence étant normale on peut intégrer terme à terme
\[ (8.2.5) \qquad \frac{1}{2\pi i} \int_{\gamma_1} \frac{f(t)\ dt}{t-z} = \sum_{n \geq 0} a_n z^n \text{ en posant :} \]
\[ (8.2.6) \qquad a_n = \frac{1}{2\pi i}\int_{\gamma_1} \frac{f(t)}{t^{n+1}}\ dt \]

\mspage{\hand{188}}
